\section{Schwarzschild 解}\label{sec:schwarzschild-solution}

真空 Einstein 场方程是非线性的, 因而远比 Laplace 方程难解. 最早发现的解之一是所谓的 \textbf{Schwarzschild 解}. 在与时间无关和球对称的简化假设下, 即寻找形如
\begin{align*}
  g&=-A^2(r)\,\mathrm{d}t\otimes\mathrm{d}t+B^2(r)\,\mathrm{d}r\otimes\mathrm{d}r\\
   &\quad+r^2\,\mathrm{d}\theta\otimes\mathrm{d}\theta
         +r^2\sin^2\theta\,\mathrm{d}\varphi\otimes\mathrm{d}\varphi
\end{align*}
的解时, 可以得到这一解, 其中 $A,B:\mathbb{R}\to\mathbb{R}$ 为未知的正值光滑函数. 注意, 当 $A\equiv B\equiv1$ 时, 这一表达式化为球坐标下的 Minkowski 度量.

容易看出, Cartan 结构方程对伪 Riemann 流形仍成立. 我们有
\[
  g=-\omega^0\otimes\omega^0+\omega^r\otimes\omega^r
    +\omega^\theta\otimes\omega^\theta+\omega^\varphi\otimes\omega^\varphi,
\]
其中
\begin{align*}
  \omega^0&=A(r)\,\mathrm{d}t;\\
  \omega^r&=B(r)\,\mathrm{d}r;\\
  \omega^\theta&=r\,\mathrm{d}\theta;\\
  \omega^\varphi&=r\sin\theta\,\mathrm{d}\varphi,
\end{align*}
故 $\{\omega^0,\omega^r,\omega^\theta,\omega^\varphi\}$ 是标准正交余标架. 第一结构方程
\begin{align*}
  \mathrm{d}\omega^\mu&=\sum_{\nu=0}^3\omega^\nu\wedge\omega_\nu^\mu;\\
  \mathrm{d}g_{\mu\nu}&=\sum_{\alpha=0}^3g_{\mu\alpha}\omega_\nu^\alpha
       +g_{\nu\alpha}\omega_\mu^\alpha,
\end{align*}
在标准正交标架下写为
\begin{align*}
  \mathrm{d}\omega^\mu&=\sum_{\nu=0}^3\omega^\nu\wedge\omega_\nu^\mu;\\
  \omega_0^0&=\omega_i^i=0;\\
  \omega_i^0&=\omega_0^i;\\
  \omega_j^i&=-\omega_i^j
\end{align*}
($i,j=1,2,3$), 连同
\begin{align*}
  \mathrm{d}\omega^0&=\frac{A'}{B}\omega^r\wedge\mathrm{d}t;\\
  \mathrm{d}\omega^r&=0;\\
  \mathrm{d}\omega^\theta&=\frac{1}{B}\omega^r\wedge\mathrm{d}\theta;\\
  \mathrm{d}\omega^\varphi&=\frac{\sin\theta}{B}\omega^r\wedge\mathrm{d}\varphi
                         +\cos\theta\,\omega^\theta\wedge\mathrm{d}\varphi,
\end{align*}
给出非零联络形式
\begin{align*}
  \omega_r^0&=\omega_0^r=\frac{A'}{B}\,\mathrm{d}t;\\
  \omega_r^\theta&=-\omega_\theta^r=\frac{1}{B}\,\mathrm{d}\theta;\\
  \omega_r^\varphi&=-\omega_\varphi^r=\frac{\sin\theta}{B}\,\mathrm{d}\varphi;\\
  \omega_\theta^\varphi&=-\omega_\varphi^\theta=\cos\theta\,\mathrm{d}\varphi.
\end{align*}
由第二结构方程
\[
  \Omega_\nu^\mu=\mathrm{d}\omega_\nu^\mu+
                \sum_{\alpha=0}^3\omega_\alpha^\mu\wedge\omega_\nu^\alpha
\]
可求得曲率形式为
\begin{align*}
  \Omega_r^0&=\Omega_0^r
    =\frac{A''B-A'B'}{AB^3}\,\omega^r\wedge\omega^0;\\
  \Omega_\theta^0&=\Omega_0^\theta
    =\frac{A'}{rAB^2}\,\omega^\theta\wedge\omega^0;\\
  \Omega_\varphi^0&=\Omega_0^\varphi
    =\frac{A'}{rAB^2}\,\omega^\varphi\wedge\omega^0;\\
  \Omega_r^\theta&=-\Omega_\theta^r
    =\frac{B'}{rB^3}\,\omega^\theta\wedge\omega^r;\\
  \Omega_r^\varphi&=-\Omega_\varphi^r
    =\frac{B'}{rB^3}\,\omega^\varphi\wedge\omega^r;\\
  \Omega_\theta^\varphi&=-\Omega_\varphi^\theta
    =\frac{B^2-1}{r^2B^2}\,\omega^\varphi\wedge\omega^\theta.
\end{align*}
因此, 由
\[
  \Omega_\nu^\mu=\sum_{\alpha<\beta}R_{\alpha\beta\nu}^{\ \ \ \mu}\,
                             \omega^\alpha\wedge\omega^\beta
\]
可从曲率形式中读出曲率张量在标准正交标架上的分量, 并进而计算 Ricci 曲率张量 $\operatorname{Ric}$ 在同一标架上的分量. $\operatorname{Ric}$ 在此标架上的非零分量为
\begin{align*}
  R_{00}&=\frac{A''B-A'B'}{AB^3}+\frac{2A'}{rAB^2};\\
  R_{rr}&=-\frac{A''B-A'B'}{AB^3}+\frac{2B'}{rB^3};\\
  R_{\theta\theta}&=R_{\varphi\varphi}
    =-\frac{A'}{rAB^2}+\frac{B'}{rB^3}+\frac{B^2-1}{r^2B^2}.
\end{align*}
因此, 真空 Einstein 场方程 $\operatorname{Ric}=0$ 等价于常微分方程组
\[
  \begin{cases}
    \dfrac{A''}{A}-\dfrac{A'B'}{AB}+\dfrac{2A'}{rA}=0\\[.7ex]
    \dfrac{A''}{A}-\dfrac{A'B'}{AB}-\dfrac{2B'}{rB}=0\\[.7ex]
    \dfrac{A'}{A}-\dfrac{B'}{B}-\dfrac{B^2-1}{r}=0
  \end{cases}
  \iff
  \begin{cases}
    \dfrac{A'}{A}+\dfrac{B'}{B}=0\\[.7ex]
    \displaystyle\left(\frac{A'}{A}\right)'+2\left(\frac{A'}{A}\right)^2
     +\frac{2A'}{rA}=0\\[.7ex]
    \dfrac{2B'}{B}+\dfrac{B^2-1}{r}=0
  \end{cases}.
\]
最后一个方程可以立即求解, 得到
\[
  B=\left(1-\frac{2m}{r}\right)^{-1/2},
\]
其中 $m\in\mathbb{R}$ 为积分常数. 第一个方程蕴含 $A=\frac{\alpha}{B}$, 其中 $\alpha>0$ 为常数. 通过重新选取时间坐标 $t$ 的尺度, 可以假定 $\alpha=1$. 最后, 容易验证第二个常微分方程恒成立. 因此, 存在所求形式的真空 Einstein 场方程的单参数解族, 由
\begin{equation}\label{eq:schwarzschild-metric}
  \begin{aligned}
    g&=-\left(1-\frac{2m}{r}\right)\mathrm{d}t\otimes\mathrm{d}t
       +\left(1-\frac{2m}{r}\right)^{-1}\mathrm{d}r\otimes\mathrm{d}r\\
     &\quad+r^2\,\mathrm{d}\theta\otimes\mathrm{d}\theta
       +r^2\sin^2\theta\,\mathrm{d}\varphi\otimes\mathrm{d}\varphi
  \end{aligned}.
\end{equation}
给出.

为解释这一解族, 计算\textbf{静止观察者}的固有加速度 [参见 Exercise~\ref{ex:special-relativity}(\ref{ex:proper-acceleration-and-spaceship})]. 这些观察者的运动是 $\frac{\partial}{\partial t}$ 的积分曲线. 若 $\{E_0,E_r,E_\theta,E_\varphi\}$ 是将 $\{\frac{\partial}{\partial t},\frac{\partial}{\partial r},\frac{\partial}{\partial\theta},\frac{\partial}{\partial\varphi}\}$ 归一化所得的标准正交标架 (因而与 $\{\omega^0,\omega^r,\omega^\theta,\omega^\varphi\}$ 对偶), 则有
\begin{align*}
  \nabla_{E_0}E_0
   &=\sum_{\mu=0}^3\omega_0^\mu(E_0)E_\mu=\omega_0^r(E_0)E_r\\
   &=\frac{A'}{AB}\omega^0(E_0)E_r
     =\frac{m}{r^2}\left(1-\frac{2m}{r}\right)^{-1/2}E_r.
\end{align*}
因此, 每个静止观察者都有一个背离原点的固有加速度
\[
  G(r)=\frac{m}{r^2}\left(1-\frac{2m}{r}\right)^{-1/2},
\]
以免向原点下落. 换言之, 他们感受到一个强度为 $G(r)$, 方向指向原点的引力场. 当 $r$ 很大时, $G(r)$ 接近由质量为 $m$ 的质点产生的 Newton 引力场中熟悉的加速度 $m/r^2$, 因此将 Schwarzschild 解解释为质量为 $m$ 的质点的广义相对论引力场. 相应地, 下文假定 $m>0$ (注意 $m=0$ 对应于 Minkowski 时空).

在求 Schwarzschild 解时, 假定 $A(r)>0$, 因而 $r>2m$. 不过, 容易验证 \eqref{eq:schwarzschild-metric} 在 $r<2m$ 时也是真空 Einstein 场方程的解. 注意, 坐标系 $(t,r,\theta,\varphi)$ 在 $r=2m$ 处是奇异的, 因此只覆盖两个不相连的开集 $\{r>2m\}$ 和 $\{r<2m\}$. 这两个开集都不是测地完备的, 例如当径向类时测地线或零模测地线趋近 $r=0$ 或 $r=2m$ 时, 就不能继续延拓 [参见 Exercise~\ref{ex:schwarzschild-solution}(\ref{ex:radial-geodesics-in-schwarzschild})]. 对于 $r=0$, 由于沿趋近这一极限的测地线的曲率会发散, 这在预料之中; 但 $r=2m$ 的情形并非如此. 事实上, 可以把这两个开集拼合起来, 得到在 $r=2m$ 处正则的真空 Einstein 场方程解. 为此, 引入所谓的 \textbf{Painlevé 时间坐标}
\[
  t'=t+\int\sqrt{\frac{2m}{r}}\left(1-\frac{2m}{r}\right)^{-1}\,\mathrm{d}r.
\]
在坐标系 $(t',r,\theta,\varphi)$ 下, Schwarzschild 度量写为
\begin{align*}
  g&=-\mathrm{d}t'\otimes\mathrm{d}t'
       +\left(\mathrm{d}r+\sqrt{\frac{2m}{r}}\mathrm{d}t'\right)
             \otimes\left(\mathrm{d}r+\sqrt{\frac{2m}{r}}\mathrm{d}t'\right)\\
   &\quad+r^2\,\mathrm{d}\theta\otimes\mathrm{d}\theta
        +r^2\sin^2\theta\,\mathrm{d}\varphi\otimes\mathrm{d}\varphi.
\end{align*}
这个表达式在 $r=2m$ 处非奇异, 并且在 $\{r>2m\}$ 和 $\{r<2m\}$ 上都是真空 Einstein 场方程的解. 由连续性, 在 $r=2m$ 处它也必为解.

子流形 $r=2m$ 称为\textbf{事件视界}, 它由零模测地线织成. 容易看出这一点, 因为 $\frac{\partial}{\partial t'}=\frac{\partial}{\partial t}$ 在 $r=2m$ 处成为零模向量, 因而其积分曲线为 (重新参数化的) 零模测地线.

考察\textbf{光锥}, 即每一点的零模切向量所成的集合, 最容易理解 Schwarzschild 时空的因果性质. 例如, 径向零模向量 $v=v^0\frac{\partial}{\partial t'}+v^r\frac{\partial}{\partial r}$ 满足
\[
  -(v^0)^2+\left(v^r+\sqrt{\frac{2m}{r}}v^0\right)^2=0
    \iff v^r=\left(\pm1-\sqrt{\frac{2m}{r}}\right)v^0.
\]
对于 $r\gg2m$, 得到近似于 Minkowski 时空的通常光锥. 然而, 随着 $r$ 趋近 $2m$, 光锥向原点\textquotedblleft 倾倒\textquotedblright, 并在 $r=2m$ 处与事件视界相切 (参见 Fig.~\ref{fig:light-cones-in-painleve-coordinates}). 由于类时曲线的切向量必须位于光锥内部, 可见穿过事件视界的粒子都不可能离开区域 $r=2m$ (正因如此, 它称为\phantomsection\label{idx:schwarzschild-black-hole}\textbf{黑洞}). 进入黑洞后, 光锥倾斜得更厉害, 迫使粒子进入奇点 $r=0$.
\RGfigure[.57\textwidth]{light-cones-in-painleve-coordinates}{Painlevé 坐标下的光锥}

注意, Painlevé 坐标下的 Schwarzschild 解在事件视界处仍不是测地完备的, 因为径向类时和零模测地线在趋近 $r=2m$ 时向过去不完备 [参见 Exercise~\ref{ex:schwarzschild-solution}(\ref{ex:radial-geodesics-in-schwarzschild})]. 在物理上, 这一点并不重要: 人们认为黑洞由 (近似) 球形的恒星坍缩而成, 恒星表面在 Fig.~\ref{fig:light-cones-in-painleve-coordinates} 的时空图中沿一条径向类时曲线运动. 由于只有恒星外部是真空, 预期 Schwarzschild 解只在该曲线上方成立, 从而去掉了 $r=2m$ 上导致不完备的区域. 尽管如此, 仍可将 Painlevé 坐标下 Schwarzschild 时空的两个副本粘合起来, 得到真空 Einstein 场方程的一个解, 它只在 $r=0$ 的两个副本处不测地完备. 这一解称为 \textbf{Kruskal 延拓}, 包含一个黑洞和它的时间反演版本, 即\textbf{白洞}.

一度有人认为, $r=0$ 处的曲率奇点是 Schwarzschild 时空高度对称性造成的假象, 更真实的恒星坍缩模型将没有奇点. Hawking 和 Penrose 证明事实并非如此: 一旦坍缩开始, 无论多么不对称, 都无法阻止奇点的形成 (参见第~\RGref{sec:hawking-singularity-theorem}~节和第~\RGref{sec:penrose-singularity-theorem}~节).

\begin{RiemannExercise}\label{ex:schwarzschild-solution}
\begin{enumerate}[label=(\arabic*),ref=\arabic*,leftmargin=*]
\item\label{ex:timelike-geodesic-triangles}
设 $(M,g)$ 为二维 Lorentz 流形.
\begin{enumerate}[label=(\alph*),leftmargin=*]
\item 考虑开集 $U\subset M$ 上的标准正交标架 $\{E_0,E_1\}$, 相应余标架为 $\{\omega^0,\omega^1\}$. 验证 Cartan 结构方程为
\begin{align*}
  \omega_1^0&=\omega_0^1;\\
  \mathrm{d}\omega^0&=\omega^1\wedge\omega_1^0;\\
  \mathrm{d}\omega^1&=\omega^0\wedge\omega_1^0;\\
  \Omega_1^0&=\mathrm{d}\omega_1^0.
\end{align*}
\item 设 $\{F_0,F_1\}$ 为另一标准正交标架, 使 $F_0\in C(E_0)$, 相应余标架为 $\{\bar\omega^0,\bar\omega^1\}$, 联络形式为 $\bar\omega_1^0$. 证明在局部有 $\sigma=\mathrm{d}u$, 其中 $\sigma=\bar\omega_1^0-\omega_1^0$, \phantomsection\label{idx:hyperbolic-angle-between-orthonormal-frames}$u$ 为 $F_0$ 与 $E_0$ 之间的双曲角 [参见 Exercise~\ref{ex:special-relativity}(\ref{ex:hyperbolic-angle})].
\item 考虑三角形 $\Delta\subset U$, 其各边为类时测地线; 设 $\alpha,\beta,\gamma$ 为各边之间的双曲角 (参见 Fig.~\ref{fig:timelike-geodesic-triangle}). 证明
\[
  \gamma=\alpha+\beta-\int_\Delta\Omega_1^0,
\]
其中按时空图的通常惯例, 将 $U$ 定向为使 $\{E_0,E_1\}$ 为\textbf{负定向}.
\item 当 $(M,g)$ 是通过固定 $(\theta,\varphi)$ 得到的 Schwarzschild 时空的全测地子流形时, 给出上述公式的物理解释 [参见第~\ref{chap:curvature}~章的 Exercise~\ref{ex:isometric-immersions}(\ref{ex:totally-geodesic-submanifolds})].
\end{enumerate}
\item\label{ex:equatorial-circular-geodesics-in-schwarzschild}
考虑具有局部坐标 $(t,r,\theta,\varphi)$ 的 Schwarzschild 时空. \textbf{赤道圆形曲线}是在这些坐标下由 $(t(\tau),r(\tau),\theta(\tau),\varphi(\tau))$ 给出的曲线, 其中 $\dot r(\tau)\equiv0$, $\theta(\tau)\equiv\frac{\pi}{2}$.
\begin{enumerate}[label=(\alph*),leftmargin=*]
\item 证明这种曲线是以固有时参数化的类时测地线的条件为
\[
  \begin{cases}
    \ddot t=0\\
    \ddot\varphi=0\\
    r\dot\varphi^2=\dfrac{m}{r^2}\dot t^2\\[.5ex]
    \left(1-\dfrac{3m}{r}\right)\dot t^2=1
  \end{cases}.
\]
由此推出, 对任意 $r>3m$, 有质量粒子都可沿圆形轨道绕中心质量运动.
\item 证明对 $r=3m$ 存在一条赤道圆形零模测地线. 位于 $r=3m$, $\theta=\frac{\pi}{2}$ 的静止观察者沿这条零模测地线的方向看, 会看到什么?
\item 自由下落的自旋粒子的角动量向量沿其运动平行移动, 且与其运动正交 [参见 Exercise~\ref{ex:general-relativity}(\ref{ex:fermi-walker-transport})]. 考虑沿圆形轨道绕点状质量 $m$ 运动的自旋粒子. 证明若角动量向量初始沿径向, 则绕行一周后它进动的角度为
\[
  \delta=2\pi\left(1-\left(1-\frac{3m}{r}\right)^{1/2}\right).
\]
(\textbf{注}: Gravity Probe B 实验期间, 人们在绕地球运行的自旋石英球上观测到了上述进动, 称为\textbf{测地进动}.)
\end{enumerate}
\item\label{ex:gravitational-redshift-in-schwarzschild}
(\emph{引力红移}) 再次考虑具有局部坐标 $(t,r,\theta,\varphi)$ 的 Schwarzschild 时空.
\begin{enumerate}[label=(\alph*),leftmargin=*]
\item 证明静止观察者在其历程中的两个事件之间测得的固有时间隔 $\Delta\tau$ 为
\[
  \Delta\tau=\left(1-\frac{2m}{r}\right)^{1/2}\Delta t,
\]
其中 $\Delta t$ 是两个事件的时间坐标之差. (\textbf{注}: 这一效应已由实验测得; 粗略地说, 引力使时间变慢.)
\item 证明若 $(t(s),r(s),\theta(s),\varphi(s))$ 是测地线, 则对任意 $\Delta t\in\mathbb{R}$, $(t(s)+\Delta t,r(s),\theta(s),\varphi(s))$ 也是测地线.
\item 利用 Fig.~\ref{fig:schwarzschild-gravitational-redshift} 中的时空图证明, 若位于 $r=r_0$ 的静止观察者测得一个光信号的周期为 $T$, 则位于 $r=r_1$ 的静止观察者对同一信号测得的周期为
\[
  T'=T\sqrt{\frac{1-\frac{2m}{r_1}}{1-\frac{2m}{r_0}}}.
\]
(\textbf{注}: 这一引力红移已由实验测得, 证实时空必有弯曲 --- 在 Minkowski 时空中必有 $T=T'$.)
\end{enumerate}
\item\label{ex:spatial-geometry-of-schwarzschild-observers}
设 $(M,g)$ 为 Schwarzschild 解的 $r>2m$ 区域, 配备 Schwarzschild 度量. $M$ 中所有静止观察者的集合是一个三维光滑流形 $\Sigma$, 具有局部坐标 $(r,\theta,\varphi)$, 且存在自然投影 $\pi:M\to\Sigma$. 如下在 $\Sigma$ 上引入 Riemann 度量 $h$: 若 $v,w\in T_{\pi(p)}\Sigma$, 则
\[
  h(v,w)=g(v^\dagger,w^\dagger),
\]
其中对每个 $u\in T_{\pi(p)}\Sigma$, 向量 $u^\dagger\in T_pU$ 满足
\[
  (\mathrm{d}\pi)_p u^\dagger=u
  \qquad\text{且}\qquad
  g\left(u^\dagger,\left(\frac{\partial}{\partial t}\right)_p\right)=0
\]
[参见 Exercise~\ref{ex:general-relativity}(\ref{ex:rotating-frame-in-minkowski-spacetime})].
\begin{enumerate}[label=(\alph*),leftmargin=*]
\item 证明 $h$ 良定义, 且
\[
  h=\left(1-\frac{2m}{r}\right)^{-1}\mathrm{d}r\otimes\mathrm{d}r
    +r^2\,\mathrm{d}\theta\otimes\mathrm{d}\theta
    +r^2\sin^2\theta\,\mathrm{d}\varphi\otimes\mathrm{d}\varphi.
\]
\item 证明 $h$ 不平坦, 但标量曲率为零.
\item 证明赤道平面 $\theta=\frac{\pi}{2}$ 与曲线 $z(r)=\sqrt{8m(r-2m)}$ 绕 $z$ 轴旋转所生成的旋转曲面等距 (参见 Fig.~\ref{fig:surface-of-revolution-of-schwarzschild-equator}).

(\textbf{注}: 这是静止观察者之间作局部距离测量所得的度量; 粗略地说, 引力使空间变形.)
\end{enumerate}
\item\label{ex:schwarzschild-orbits-and-light-deflection}
本题详细研究 Schwarzschild 时空的类时测地线和零模测地线. 首先注意, 子流形 $\theta=\frac{\pi}{2}$ 是全测地的 [参见第~\ref{chap:curvature}~章的 Exercise~\ref{ex:isometric-immersions}(\ref{ex:totally-geodesic-submanifolds})]. 适当选取角坐标 $(\theta,\varphi)$, 总可以假设测地线的初始条件与这一子流形相切; 因而只需研究三维 Lorentz 流形 $(M,g)$ 的类时及零模测地线, 其中
\[
  g=-\left(1-\frac{2m}{r}\right)\mathrm{d}t\otimes\mathrm{d}t
    +\left(1-\frac{2m}{r}\right)^{-1}\mathrm{d}r\otimes\mathrm{d}r
    +r^2\,\mathrm{d}\varphi\otimes\mathrm{d}\varphi.
\]
\begin{enumerate}[label=(\alph*),leftmargin=*]
\item 证明 $\frac{\partial}{\partial t}$ 和 $\frac{\partial}{\partial\varphi}$ 是 Killing 向量场 [参见第~\ref{chap:riemannian-manifolds}~章的 Exercise~\ref{ex:levi-civita-connection}(\ref{ex:killing-vector-field})].
\item 由此推出, 曲线 $c:\mathbb{R}\to M$ 是指向未来的测地线 (若为类时, 则以固有时参数化) 的方程可以写为
\[
  \begin{cases}
    g(\dot c,\dot c)=-\sigma\\[.5ex]
    g\left(\dfrac{\partial}{\partial t},\dot c\right)=-E\\[.5ex]
    g\left(\dfrac{\partial}{\partial\varphi},\dot c\right)=L
  \end{cases}
  \iff
  \begin{cases}
    \dot r^2=E^2-\left(\sigma+\dfrac{L^2}{r^2}\right)\left(1-\dfrac{2m}{r}\right)\\[.6ex]
    \left(1-\dfrac{2m}{r}\right)\dot t=E\\[.6ex]
    r^2\dot\varphi=L
  \end{cases}
\]
其中 $E>0$ 和 $L$ 为积分常数, 类时测地线取 $\sigma=1$, 零模测地线取 $\sigma=0$.
\item 证明若 $L\ne0$, 则 $u=\frac{1}{r}$ 满足
\[
  \frac{\mathrm{d}^2u}{\mathrm{d}\varphi^2}+u=\frac{m\sigma}{L^2}+3mu^2.
\]
\item 若相对论修正很小, 则有 $mu\ll1$, 因此近似方程
\[
  \frac{\mathrm{d}^2u}{\mathrm{d}\varphi^2}+u=\frac{m}{L^2}
\]
对类时测地线成立. 证明该方程的解是极坐标下由下式给出的圆锥曲线:
\[
  u=\frac{m}{L^2}(1+\varepsilon\cos(\varphi-\varphi_0)),
\]
其中积分常数 $\varepsilon\geq0$ 及 $\varphi_0$ 分别为偏心率和近心点幅角.
\item 证明当 $\varepsilon\ll1$ 时, 上述近似解满足
\[
  u^2\simeq\frac{2m}{L^2}u-\frac{m^2}{L^4}.
\]
说明为何接近圆形轨道且相对论修正很小的类时测地线, 给出方程
\[
  \frac{\mathrm{d}^2u}{\mathrm{d}\varphi^2}
  +\left(1-\frac{6m^2}{L^2}\right)u
  =\frac{m}{L^2}\left(1-\frac{3m^2}{L^2}\right)
\]
的近似解, 因而近心点每周约前进 $\frac{6\pi m}{r}$ 弧度.

(\textbf{注}: 广义相对论的第一次成功就得益于这一效应, 它解释了\textbf{水星近日点的反常进动} --- 每世纪 43 角秒.)
\item 证明若忽略相对论修正, 则零模测地线满足
\[
  \frac{\mathrm{d}^2u}{\mathrm{d}\varphi^2}+u=0.
\]
证明该方程的解是极坐标下直线的方程
\[
  u=\frac{1}{b}\sin(\varphi-\varphi_0),
\]
其中积分常数 $b>0$ 和 $\varphi_0$ 分别为\textbf{碰撞参数} (到中心的最小距离) 以及直线与 $x$ 轴的夹角.
\item 假定 $mu\ll1$. 通过寻找形如
\[
  u=\frac{1}{b}\left(\sin\varphi+\frac{m}{b}v\right)
\]
的近似解来纳入相对论修正 (为简单起见, 取 $\varphi_0=0$). 证明 $v$ 是方程
\[
  \frac{\mathrm{d}^2v}{\mathrm{d}\varphi^2}+v=3\sin^2\varphi
\]
的近似解, 因而 $u$ 近似由
\[
  u=\frac{1}{b}\left(\sin\varphi
     +\frac{m}{b}\left(\frac{3}{2}+\frac{1}{2}\cos(2\varphi)
                              +\alpha\cos\varphi+\beta\sin\varphi\right)\right)
\]
给出, 其中 $\alpha,\beta$ 为积分常数.
\item 证明对于零模测地线的入射部分 ($\varphi\simeq0$), 近似有
\[
  u=0\iff\varphi=-\frac{m}{b}(2+\alpha).
\]
类似地, 证明对于出射部分 ($\varphi\simeq\pi$), 近似有
\[
  u=0\iff\varphi=\pi+\frac{m}{b}(2-\alpha).
\]
推出沿其路径, $\varphi$ 的改变约为
\[
  \Delta\varphi=\pi+\frac{4m}{b}
\]
弧度, 因此零模测地线向中心偏折约 $\frac{4m}{b}$ 弧度.

(\textbf{注}: 测量太阳引起的这种\textbf{光线偏折} --- $1.75$ 角秒 --- 首次在实验上证实了广义相对论, 使 Einstein 一夜之间成为世界名人.)
\end{enumerate}
\item\label{ex:birkhoff-spherical-vacuum-theorem}
(\emph{Birkhoff 定理}) 证明在局部坐标 $(t,r,\theta,\varphi)$ 下由
\[
  g=A^2(t,r)\,\mathrm{d}t\otimes\mathrm{d}t+B^2(t,r)\,\mathrm{d}r\otimes\mathrm{d}r
     +r^2\,\mathrm{d}\theta\otimes\mathrm{d}\theta
     +r^2\sin^2\theta\,\mathrm{d}\varphi\otimes\mathrm{d}\varphi
\]
给出的 Ricci 平坦 Lorentz 度量只能是 Schwarzschild 度量. 粗略地说, 球对称的质量位形不产生辐射.
\item\label{ex:radial-geodesics-in-schwarzschild}
\begin{enumerate}[label=(\alph*),leftmargin=*]
\item 证明 Schwarzschild 时空的区域 $\{r>2m\}$ 和 $\{r<2m\}$ 中的径向类时或零模测地线, 在趋近 $r=0$ 或 $r=2m$ 时不能延拓.
\item 证明 Schwarzschild 时空的 Painlevé 延拓中的径向类时或零模测地线, 在趋近 $r=2m$ 时可以向未来延拓, 但不能向过去延拓.
\item 证明在 Painlevé 坐标下满足
\[
  \frac{\mathrm{d}r}{\mathrm{d}t'}=-\sqrt{\frac{2m}{r}}
\]
的径向观察者是自由下落的, 且 $t'$ 是他们的固有时.
\item 当一个粒子落入黑洞时, 静止观察者会看到什么?
\end{enumerate}
\item\label{ex:proper-time-from-horizon-to-singularity}
证明穿过视界的观察者将在 $\Delta\tau\leq\pi m$ 的固有时间隔内到达奇点.
\end{enumerate}
\end{RiemannExercise}
\RGfigure[.37\textwidth]{timelike-geodesic-triangle}{类时测地三角形}
\RGfigure[.53\textwidth]{schwarzschild-gravitational-redshift}{引力红移}
\RGfigure[.48\textwidth]{surface-of-revolution-of-schwarzschild-equator}{与赤道平面等距的旋转曲面}
